\thispagestyle{empty}

\begin{center}
\Large
\textbf{誌\hspace{1cm}謝}
\end{center}

\vspace{1cm}
{
\linespread{2}%
\selectfont
\hspace{0.25cm}

% 下面開始寫致謝內容
能完成研究所這段旅程，首先要感謝辛苦的爸媽跟哥哥，幾乎幫我處理很多生活大小事，讓我能無後顧之憂地專心做研究以及修課，將多數精力放在學業上。

再來要感謝林甫俊教授從碩一開始的指導，每周的會議，Lab，提供的論文，以及問答都讓我獲益良多，碩二引導我將一個抽象的題目逐步轉化成一個具體的論文。
還有陳志成教授以及陳健教授，將與CNTi合作的事情指派給我讓我能夠從中學習與成長。安排許多國外教授演講分享他們的工作或者人生經驗。

實驗室的學長姐以及同屆的實驗室夥伴更是努力前進的原動力，一周七天每天去實驗室還能看到大家一起卷就不會覺得累。感謝吳國成與官劉翔一起討論所有Lab,研究相關的問題，在事情一起轟炸，心很累的晚上一起吃宵夜。
感謝曹詠喧，鄭珮綺，黃鈞偉一起聊課業有關或無關的內容，討論各種生活遇到的或有趣或很雖的瑣事。實驗室五排，大家一起吃宵夜聊天是很快樂的時光。
感謝杜峯一起解決許多任務，提供很多free5GC相關的幫助，完成一個月極限的國科會Demo。

感謝Ian學長，帶領我們入門free5GC，花時間與我們討論issue解法，在解決CNTi issue的過程中提供很多建議，還會分享很多有趣的專案。
也感謝學長姐們(威成，澤田，定遠，Michael，BC，品凡，昆霖，與234實驗室的許多人)提供很多過去的考古或code，分享很多就業的內容，以及許多有趣的故事，一起下戰棋。另外感謝謝孟瀚無私分享所有AI工具相關的技巧，林哲偉一起通過艱難的OSC。

\vspace{3cm}
\begin{flushright}
\studentCnName \ 於

國立陽明交通大學\ \NameofDepartmentInstituteCN

\ThesisDateTW
\end{flushright}
}